\documentclass[12pt]{article}
\usepackage[margin=1in]{geometry}
\usepackage{amsmath, amssymb}
\usepackage{enumitem, array}

\usepackage{boxedminipage}
\usepackage{fancyhdr}
\usepackage{ragged2e, comment, parskip}


\newenvironment{problem}[1][]{%
  \bigskip% put space before problem statement box %
  \noindent\begin{boxedminipage}{\columnwidth}%
  \def\@tempa{#1}%
  \ifx\@tempa\empty\else%
    \textbf{#1}\hspace{0.5em}%
  \fi%
}{%
  \end{boxedminipage}%
}


\begin{document}

\begin{center}
    \Large \textbf{MATH 4790 Final Exam Spring 2026} \\[1em]
\end{center}

\vspace{.25cm}

\noindent Name: \underline{\hspace{8cm}} \hspace{.25cm} Auburn ID: \underline{\hspace{4cm}}



\section*{Instructions}
\begin{enumerate}[label=\arabic*.]
\item Print your full name and your banner ID on your exam. Then finish reading these instructions, and sign the bottom of this page.
\item All electronic equipment (cell phones, Apple Watch, AirPod, etc) are \textbf{NOT ALLOWED} during the exam time.
\item You are allowed to use up to 2 calculators from the following list: BA-35, BA II Plus, BA II Plus Professional, TI-30Xa, TI-30X II (IIS solar or IIB battery), TI-30XS MultiView (or XB battery).
\item No talking is allowed during the exam.
\item Unless otherwise indicated, \textbf{SHOW ALL YOUR WORK}. If no work is shown, no partial credit can be awarded. Your work and answer need to be accurate and relevant to receive points.
\item You will be given exactly 120 minutes for this exam.
\item Any student not following the above instructions nor behaving according to the above instructions during the exam may have their exam confiscated and points deducted.
\end{enumerate}

I have read and fully understand all of the above instructions. \\[2em]
Signature: \hrulefill


\vspace{.5cm}
Do not write in the table below

% scoring table
\begin{center}

\begin{tabular}{c c}

% -------- LEFT TABLE --------
\begin{tabular}{|*{3}{>{\centering\arraybackslash}p{.1\linewidth}|}}
    \hline
    \textbf{Problem} & \textbf{Points} & \textbf{Score} \\
    \hline
    1 & 5 &  \\
    \hline
    2 & 5 &  \\
    \hline
    3 & 5 &  \\
    \hline
    4 & 5 &  \\
    \hline
    5 & 5 &  \\
    \hline
    6 & 5 &  \\
    \hline
    7 & 5 &  \\
    \hline
    8 & 5 &  \\
    \hline
    9 & 5 &  \\
    \hline
    10 & 5 &  \\
    \hline
\end{tabular}

&

% -------- RIGHT TABLE --------
\begin{tabular}{|*{3}{>{\centering\arraybackslash}p{.1\linewidth}|}}
    \hline
    \textbf{Problem} & \textbf{Points} & \textbf{Score} \\
    \hline
    11 & 5 &  \\
    \hline
    12 & 5 &  \\
    \hline
    13 & 5 &  \\
    \hline
    14 & 5 &  \\
    \hline
    15 & 5 &  \\
    \hline
    16 & 5 &  \\
    \hline
    17 & 5 &  \\
    \hline
    18 & 5 &  \\
    \hline
    19 & 5 &  \\
    \hline
    20 & 5 &  \\
    \hline
\end{tabular}

\end{tabular}

\vspace{1em}

% -------- TOTAL ROW --------
\begin{tabular}{|c|c|c|}
    \hline
    \textbf{Total} & 100 & \text{\hspace{2cm}}\\
    \hline
\end{tabular}

\end{center}

\newpage

\newpage

\begin{center}
{\Large \textbf{Answer Sheet}}
\end{center}

\vspace{1cm}

\begin{center}
\begin{minipage}{0.4\textwidth}
\begin{enumerate}
\item \underline{\hspace{1.5cm}}

\vspace{0.6cm}

\item \underline{\hspace{1.5cm}}

\vspace{0.6cm}

\item \underline{\hspace{1.5cm}}

\vspace{0.6cm}

\item \underline{\hspace{1.5cm}}

\vspace{0.6cm}

\item \underline{\hspace{1.5cm}}

\vspace{0.6cm}

\item \underline{\hspace{1.5cm}}

\vspace{0.6cm}

\item \underline{\hspace{1.5cm}}

\vspace{0.6cm}

\item \underline{\hspace{1.5cm}}

\vspace{0.6cm}

\item \underline{\hspace{1.5cm}}

\vspace{0.6cm}

\item \underline{\hspace{1.5cm}}
\end{enumerate}
\end{minipage}
\hspace{0.15\textwidth}
\begin{minipage}{0.4\textwidth}
\begin{enumerate}
\setcounter{enumi}{10}
\item \underline{\hspace{1.5cm}}

\vspace{0.6cm}

\item \underline{\hspace{1.5cm}}

\vspace{0.6cm}

\item \underline{\hspace{1.5cm}}

\vspace{0.6cm}

\item \underline{\hspace{1.5cm}}

\vspace{0.6cm}

\item \underline{\hspace{1.5cm}}

\vspace{0.6cm}

\item \underline{\hspace{1.5cm}}

\vspace{0.6cm}

\item \underline{\hspace{1.5cm}}

\vspace{0.6cm}

\item \underline{\hspace{1.5cm}}

\vspace{0.6cm}

\item \underline{\hspace{1.5cm}}

\vspace{0.6cm}

\item \underline{\hspace{1.5cm}}
\end{enumerate}
\end{minipage}
\end{center}

\newpage

%%%%% PROBLEM 1 %%%%%
\newpage
\noindent \textbf{Problem 1.}
An investor deposits $Z$ into a fund that accumulates at an annual effective interest rate of $4\%$ for 12 years.

At the same time, $\frac{Z}{2}$ is invested in a second fund that accumulates at an annual effective rate of discount $d$ for 9 years.

The total interest earned in each fund over its respective investment period is the same.

Determine $d$.

\vspace{.5cm}

\begin{enumerate}[label = \Alph*.]
\item $4.8\%$
\item $5.6\%$
\item $6.2\%$
\item $6.9\%$
\item $7.5\%$
\end{enumerate}

\newpage



\noindent \textbf{Problem 2.}
A deposit of \$250 is made into an account at time 0. The force of interest, $\delta_t$, is given by

\[
\delta_t =
\begin{cases}
0.035, & 0 < t \leq 3,\\
0.002t^2, & 3 < t \leq 6,\\
\dfrac{1}{t+20}, & 6 < t \leq 10.
\end{cases}
\]

Determine the accumulated value of the deposit at time 10.

\vspace{.5cm}

\begin{enumerate}[label = \Alph*.]
\item \$341.16
\item \$352.80
\item \$363.42
\item \$375.09
\item \$389.64
\end{enumerate}

\newpage


\noindent \textbf{Problem 3.}
At an effective annual interest rate $i$, you are given:

\begin{enumerate}[label = \roman*.]
\item the present value of an annuity-immediate with annual payments of 1 for $n$ years is 20;
\item the present value of an annuity-immediate with annual payments of 1 for $2n$ years is 32.
\end{enumerate}

Calculate the present value of an annuity-due with annual payments of 1 for $3n$ years.

\vspace{.5cm}

\begin{enumerate}[label = \Alph*.]
\item 40
\item 42
\item 44
\item 46
\item 48
\end{enumerate}

\newpage


\noindent \textbf{Problem 4.}
An endowment is valued using an annual effective interest rate of $4\%$. The endowment provides annual payments as follows:

\begin{enumerate}[label = \roman*.]
\item The first 12 payments are each \$1{,}800, with the first payment occurring two years from today.
\item Starting with the 13th payment, each payment is $k\%$ larger than the previous payment.
\item The present value of all payments is \$76{,}977.
\end{enumerate}

Calculate $k$.

\vspace{.5cm}

\begin{enumerate}[label = \Alph*.]
\item $1.75$
\item $2.00$
\item $2.25$
\item $2.50$
\item $2.75$
\end{enumerate}

\newpage

\noindent \textbf{Problem 5.}
An investor makes deposits of \$1{,}200 at the end of each year for 6 years. Each deposit earns interest at an annual effective rate of $9\%$, but the interest received is reinvested at a constant force of interest $\delta = 0.05$.

Determine the accumulated value at the end of 6 years.

\vspace{.5cm}

\begin{enumerate}[label = \Alph*.]
\item \$8{,}850
\item \$8{,}935
\item \$9{,}020
\item \$9{,}115
\item \$9{,}240
\end{enumerate}

\newpage

\noindent \textbf{Problem 6.}
A perpetuity pays continuously. At time $t$, the payment rate is $4t$. Time is measured in years.

At an annual effective interest rate of $6\%$, calculate the present value of the perpetuity.

\vspace{.5cm}

\begin{enumerate}[label = \Alph*.]
\item \$820.45
\item \$875.32
\item \$902.78
\item \$918.27
\item \$945.10
\end{enumerate}

\newpage

\noindent \textbf{Problem 7.}
Two annuities have the same present value at an annual effective interest rate of $8\%$.

The first annuity is an annuity-immediate. Its first payment is \$1{,}000, and each payment is $3\%$ larger than the previous payment. This annuity makes payments for $n$ years.

The second annuity makes payments at the beginning of each year forever. The first payment is \$100, the second payment is \$140, the third payment is \$180, and so on.2

Determine $n$.

\vspace{.5cm}

\begin{enumerate}[label = \Alph*.]
\item 9
\item 10
\item 11
\item 12
\item 13
\end{enumerate}

\newpage

\noindent \textbf{Problem 8.}
A \$85{,}000 loan has an annual nominal interest rate of $7.2\%$ convertible monthly. The loan is repaid with monthly payments, with the first payment due one month from the date of the loan.

For the first four years, each monthly payment is \$950. All payments thereafter are \$1{,}300 except for a final smaller balloon payment.

Calculate the balloon payment.

\vspace{.5cm}

\begin{enumerate}[label = \Alph*.]
\item \$1{,}062
\item \$1{,}084
\item \$1{,}109
\item \$1{,}137
\item \$1{,}165
\end{enumerate}

\newpage


\noindent \textbf{Problem 9.}
A 30-year loan is repaid with equal annual payments.

The amount of interest paid in the 7th payment is \$120.

The amount of interest paid in the 19th payment is \$96.

Calculate the original amount of the loan.

\vspace{.5cm}

\begin{enumerate}[label = \Alph*.]
\item \$985
\item \$1{,}013
\item \$1{,}046
\item \$1{,}081
\item \$1{,}125
\end{enumerate}

\newpage

\noindent \textbf{Problem 10.}
A loan is repaid with level monthly payments at the end of each month for $n^2+n-10$ months.

The annual nominal interest rate is $21\%$, convertible monthly.

The $n$th payment consists of approximately equal amounts of interest and principal.

Calculate $n$.

\vspace{.5cm}

\begin{enumerate}[label = \Alph*.]
\item 5
\item 6
\item 7
\item 8
\item 9
\end{enumerate}




\newpage


\noindent \textbf{Problem 11.}
An $n$-year bond with annual coupons has the following characteristics:

\begin{enumerate}[label = \roman*.]
\item The redemption value at maturity is \$1{,}000;
\item The annual effective yield rate is $6\%$;
\item The book value immediately after the 12th coupon is \$748.99; and
\item The book value immediately after the 13th coupon is \$753.93.
\end{enumerate}

Calculate $n$.

\vspace{.5cm}

\begin{enumerate}[label = \Alph*.]
\item 32
\item 34
\item 36
\item 38
\item 40
\end{enumerate}

\newpage

\noindent \textbf{Problem 12.}
Bond A and Bond B are each 6-year \$1{,}000 face amount bonds. In addition:

\begin{enumerate}[label = \roman*.]
\item Bond A has an annual coupon rate of $5.6\%$ paid quarterly.
\item Bond B has an annual coupon rate of $4\%$ paid semiannually.
\item The price of Bond B is \$50 less than the price of Bond A.
\item The annual nominal yield rate for Bond A is $4.8\%$ convertible quarterly.
\end{enumerate}

Calculate the annual nominal yield rate for Bond B, convertible semiannually.

\vspace{.5cm}

\begin{enumerate}[label = \Alph*.]
\item $3.96\%$
\item $4.08\%$
\item $4.16\%$
\item $4.24\%$
\item $4.35\%$
\end{enumerate}

\newpage

\noindent \textbf{Problem 13.}
A 20-year bond with semiannual coupons has a redemption value of \$1{,}000. It is purchased at a discount to yield $8\%$ nominal, convertible semiannually.

If the amount of discount accumulated in the 35th coupon payment is \$18.50, which of the following is closest to the total discount in the original purchase price?

\vspace{.5cm}

\begin{enumerate}[label = \Alph*.]
\item \$118
\item \$120
\item \$122
\item \$124
\item \$126
\end{enumerate}

\newpage

\noindent \textbf{Problem 14.}
Baxter Investments purchases an $n$-year \$1{,}000 face amount bond with $5\%$ annual coupons at a price of \$920, to yield an annual effective rate $i$, where $i>0$.

After the first three years, the bond's book value has increased by \$46.29.

Calculate $i$.

\vspace{.5cm}

\begin{enumerate}[label = \Alph*.]
\item $6.0\%$
\item $6.5\%$
\item $7.0\%$
\item $7.5\%$
\item $8.0\%$
\end{enumerate}

\newpage

\noindent \textbf{Problem 15.}
Morgan is considering purchasing a \$1{,}000-par callable bond with semiannual coupons at an annual coupon rate of $7\%$. The bond matures at the end of 25 years and has a redemption value of \$1{,}000.

The bond may be called as follows:

\begin{itemize}
    \item at the end of years 10 through 12 for \$1{,}080;
    \item at the end of years 13 through 17 for \$1{,}050;
    \item at the end of years 18 through 24 for \$1{,}020.
\end{itemize}

Calculate the highest price Morgan should pay to guarantee a yield of at least $6\%$ nominal, convertible semiannually.

\vspace{.5cm}

\begin{enumerate}[label = \Alph*.]
\item \$1{,}106
\item \$1{,}110
\item \$1{,}114
\item \$1{,}118
\item \$1{,}122
\end{enumerate}

\newpage

\noindent \textbf{Problem 16.}
You are given the following term structure of effective annual spot rates on zero-coupon bonds:

\begin{itemize}
\item 1-year maturity: $5\%$
\item 2-year maturity: $6\%$
\item 3-year maturity: $7\%$
\end{itemize}

The price of a 4-year \$100 par value bond with annual coupons of \$7 is \$98.01.

Calculate the two-year forward rate, deferred for 2 years.

\vspace{.5cm}

\begin{enumerate}[label = \Alph*.]
\item $8.4\%$
\item $9.0\%$
\item $9.5\%$
\item $10.1\%$
\item $10.7\%$
\end{enumerate}

\newpage

\noindent \textbf{Problem 17.}
You are given the following information about Red Oak Financial.

\begin{enumerate}[label = \roman*.]
\item It has a single liability of \$1{,}500{,}000 due 11 years from today.
\item Its asset portfolio consists of a zero-coupon bond maturing in 6 years for \$522{,}351 and a zero-coupon bond maturing in $n$ years.
\item At an annual effective interest rate of $5\%$, Red Oak Financial's position is fully immunized.
\end{enumerate}

Calculate $n$.

\vspace{.5cm}

\begin{enumerate}[label = \Alph*.]
\item 13
\item 14
\item 15
\item 16
\item 17
\end{enumerate}

\newpage

\noindent \textbf{Problem 18.}
A perpetuity-immediate pays \$50 at the end of each year forever.

Using an annual effective interest rate of $8\%$, let
\[
Z = \frac{\text{Macaulay Duration}+\text{Modified Duration}}{2}
\]
denote the average of the Macaulay and modified durations.

Calculate $Z$.

\vspace{.5cm}

\begin{enumerate}[label = \Alph*.]
\item 12.50
\item 13.50
\item 14.50
\item 15.50
\item 16.50
\end{enumerate}

\newpage

\noindent \textbf{Problem 19.}
A company has liabilities of \$1{,}000 due in 1 year and \$2{,}000 due in 2 years.

The company cash flow exactly matches these liabilities using the following assets:

\begin{itemize}
\item a 2-year bond with annual coupons of $5\%$ yielding $5\%$
\item a 1-year bond with annual coupons of $4\%$ yielding $4\%$
\end{itemize}

Calculate the internal rate of return (IRR) on the asset portfolio.

\vspace{.5cm}

\begin{enumerate}[label = \Alph*.]
\item $4.00\%$
\item $4.25\%$
\item $4.50\%$
\item $4.75\%$
\item $5.00\%$
\end{enumerate}

\newpage

\noindent \textbf{Problem 20.}
A portfolio has the following cash flows:

\begin{itemize}
\item \$105 at time 1;
\item \$220.50 at time 2;
\item \$378.69 at time 4.
\end{itemize}

The annual effective interest rate is $5\%$ for the first 2 years. After time 2, the annual effective interest rate is $7\%$.

Calculate the Macaulay convexity of the portfolio.

\vspace{.5cm}

\begin{enumerate}[label = \Alph*.]
\item $8.50$
\item $9.00$
\item $9.50$
\item $10.00$
\item $10.50$
\end{enumerate}

\newpage


Scratch Work




\end{document}
